% vim: ai sw=2

\documentclass[handout,xcolor=dvipsnames]{beamer}

\let\Oldemph\emph
\renewcommand{\emph}[1]{\textcolor{blue}{\Oldemph{#1}}}
\newcommand{\heading}[1]{\textbf{\textcolor{BurntOrange}{#1}}}

% Copyright 2004 by Till Tantau <tantau@users.sourceforge.net>.
%
% In principle, this file can be redistributed and/or modified under
% the terms of the GNU Public License, version 2.
%
% However, this file is supposed to be a template to be modified
% for your own needs. For this reason, if you use this file as a
% template and not specifically distribute it as part of a another
% package/program, I grant the extra permission to freely copy and
% modify this file as you see fit and even to delete this copyright
% notice. 


\mode<presentation>
{
%  \usetheme{Warsaw}
\usetheme{Madrid}

  \setbeamercovered{transparent}
}

\usepackage[utf8]{inputenc}
\usepackage{times}
\usepackage[T1]{fontenc}

\usepackage{graphicx}
\usepackage[final]{listings}
\usepackage{multirow}

\lstset{tabsize=2,numbers=none,showstringspaces=false,breaklines=true,
	breakatwhitespace=true,escapeinside={(*@}{@*)},
	basicstyle=\scriptsize,keywordstyle=\bfseries\color{OliveGreen},
	commentstyle=\color{blue}}

\title[Stanse] % (optional, use only with long paper titles)
{Výlet do jádra Linuxu}

% \subtitle
% {Include Only If Paper Has a Subtitle}

\author[J. Slabý]{\emph{Jiří~Slabý}}

\institute[FI] % (optional, but mostly needed)
{Fakulta Informatiky, Masarykova Univerzita}

\date[17. 5. 2010, Brno]{Brno, 17. května 2010}

\subject{Theoretical Computer Science}

\begin{document}

\begin{frame}
  \titlepage
\end{frame}

%%%%%%%%%%%%%%%%%%%%%%%

\begin{frame}
  \frametitle{Osnova přednášky}
  \begin{itemize}
  \item Úvod do světa jádra
  \item Rozdíly uživatelské programy vs. jádro
  \item Základy programování jádra
  \item Příklady (alokace, komunikace s procesy)
  \item Závěr
  \item Navazuje: přednáška R. Krejčího
  \end{itemize}
\end{frame}

\begin{frame}
  \frametitle{Úvod}
  \begin{itemize}
  \item Co je to jádro OS?
  \begin{itemize}
  \item Prostředník mezi procesy a HW
  \item Ovládá a zpřístupňuje jej
  \end{itemize}
  \item Běží ve speciálním režimu
  \item Programování vyžaduje specifické znalosti (jako každé jiné)
  \end{itemize}

  \pause
  \medskip

  \heading{Proč umět programovat jádro?}
  \begin{itemize}
  \item Nový hardware (žádné ovladače)
  \item Starý hardware (špatné ovladače)
  \item Nová funkcionalita (lepší plánovač, síťový stack, atd.)
  \item Oprava chyb
  \end{itemize}

\end{frame}

\begin{frame}
  \frametitle{Rozdíly program vs. jádro}
  \heading{Program (v C)}
  \begin{itemize}
  \item Knihovny
  \begin{itemize}
  \item \emph{libc} (standardní funkce typu \texttt{printf}, \texttt{strlen},
    \texttt{memset}, \ldots)
  \item Ostatní -- přidávající funckionalitu (\emph{boost}, \emph{pthread},
    \ldots)
  \end{itemize}
  \item Oddělený paměťový prostor
  \begin{itemize}
  \item Procesy se neovlivňují, pakliže to nebylo explicitně povoleno -- zápis
    na adresu mimo alokovaných/mapovaných zabije aplikaci.
  \end{itemize}
  \item Počáteční funkce \texttt{main}
  \end{itemize}

  \pause
  \medskip

  \heading{Jádro}
  \begin{itemize}
  \item Nic z výše uvedeného
  \item \emph{Chyba v jádře často působí pád systému}
  \item Některé funkce reimplementovány
  \begin{itemize}
  \item \texttt{printf} jako \texttt{printk} v Linuxu, \texttt{printf} v
    BSD
  \end{itemize}
  \item Žádná aritmetika s pohyblivou desetinnou čárkou apod.
  \end{itemize}
\end{frame}

\begin{frame}
  \frametitle{Linuxové jádro}
  \heading{Kód}
  \begin{itemize}
  \item GNU C
  \begin{itemize}
  \item \texttt{gcc} alespoň 3.4.2
  \end{itemize}
  \item Modulární
  \item Víceméně pevně daný styl
  \begin{itemize}
  \item Tzv. \emph{CodingStyle}
  \end{itemize}
  \item Neobsahuje boot-strap
  \begin{itemize}
  \item Práce pro GRUB, (E)LILO, QEMU, apod.
  \item Provedou zavedení do paměti a skok do funkce \texttt{start\_kernel}
  \end{itemize}
  \item Linkuje se do jednoho celku -- monolit
  \begin{itemize}
  \item Volání přímo (žádné předávání zpráv ve smyslu mikrojádra)
  \item Avšak dovoluje přidávat (přilinkovat) za běhu kód ve formě modulů
  \end{itemize}
  \item Počáteční funkci ("main") souboru definuje makro \texttt{module\_init}
  \item Příklad
  \end{itemize}

\end{frame}

\begin{frame}[fragile]
  \frametitle{Modul "Ahoj světe"}
  \begin{center}
  \begin{tabular}{|c|c|} \hline
  code.c & Makefile \\ \hline
  \begin{lstlisting}[language=C]
#include <linux/module.h>

static int my_init(void)
{
  printk("Ahoj svete?\n");
  return 0;
}

static void my_exit(void)
{
  printk("Ahoj svete!\n");
}

module_init(my_init);
module_exit(my_exit);
  \end{lstlisting} &
  \begin{lstlisting}[language=XML]
KDIR=/lib/modules/$(shell uname -r)/build
KBUILD=$(MAKE) -C $(KDIR) M=$(PWD)

obj-m := code.o

default:
  $(KBUILD) modules
  \end{lstlisting} \\ \hline
  \end{tabular}

  \pause

  \begin{tabular}{|c|c|}
  \begin{lstlisting}[language=XML]
# insmod code.ko
# dmesg|tail -1
Ahoj svete?
  \end{lstlisting} &
  \begin{lstlisting}[language=XML]
# rmmod code
# dmesg|tail -1
Ahoj svete!
  \end{lstlisting}
  \end{tabular}
  \end{center}
\end{frame}

\begin{frame}
  \frametitle{Linuxové jádro}

  \heading{Chování}
  \begin{itemize}
  \item Velká část je řízena událostmi (event-driven)
  \begin{itemize}
  \item Přerušení od HW, požadavky uživatele, komunikace modulů, atd.
  \item Implicitní souběžnost kódu
  \end{itemize}
  \item 2 úrovně zabezpečení
  \begin{itemize}
  \item \emph{Uživatelský prostor} -- kontrolované vykonávání instrukcí
  \item \emph{Prostor jádra} -- neomezené vykonávání (kód může cokoliv)
  \item Na x86 odpovídá privilege ring 0 a 3
  \end{itemize}
  \item Několik kontextů vykonávání kódu
  \begin{itemize}
  \item \emph{Proces} -- proces zavolá jádro (čtení souboru)
  \item \emph{Prerušení} -- kontext právě přerušeného procesu (síť hlásí paket)
  \item \emph{Měkké přerušení} -- taktéž, ale prerušení jsou povolena (časovač)
  \end{itemize}
  \end{itemize}

\end{frame}

\begin{frame}
  \frametitle{Diskuze}
  \begin{center} \huge
  Děkuji za pozornost! \\[1.8em]
  Dotazy?
  \end{center}

  \pause
  \bigskip

  \begin{center} \large
  Přivítejme Radka Krejčího s navazující přednáškou.
  \end{center}
\end{frame}

\end{document}

